% GENERATED FILE. Edit canonical lessons, then run npm run generate:books.
% canonical-source-hash: fnv1a64:274a9ca597362f64
% canonical-lessons: PA-C241-chhillna, PA-C241-katarna, PA-C241-chirna, PA-C241-mirach, PA-C241-vesan

\chapter{Cutting Vegetables}
\label{ch:pa-a2-cutting-vegetables}
\hlchaptermodality{241}

\begin{chapteropening}
\textbf{By the end of this chapter:} \emph{I can say how I peel and cut vegetables.}

Complete the last lesson of chapter 241: I can say how I peel and cut vegetables.
\end{chapteropening}

\section[\texorpdfstring{\pa{ਛਿੱਲਣਾ} (chhillṇā)}{ਛਿੱਲਣਾ (chhillṇā)}]{\pa{ਛਿੱਲਣਾ} (chhillṇā) — to peel}

\label{lesson:PA-C241-chhillna}



\begin{quote}
Before the new one: say the Punjabi for a pestle, then the Punjabi for to pound.
\end{quote}

\subsection*{You'll want to know: \pa{ਛਿੱਲਣਾ}}
\textbf{\pa{ਛਿੱਲਣਾ}} — \emph{chhillṇā} — \textquotedblleft{}to peel\textquotedblright{}.

\begin{quote}
\pa{ਮੈਂ} \pa{ਚਾਕੂ} \pa{ਨਾਲ} \pa{ਆਲੂ} \pa{ਛਿੱਲਦੀ} \pa{ਹਾਂ।}
\end{quote}

\begin{quote}
\emph{maiṁ chākū nāl ālū chhilldī hāṁ.}
\end{quote}

\begin{quote}
I peel the potatoes with a knife.
\end{quote}

Past (I …): \pa{ਮੈਂ} \pa{ਆਲੂ} \pa{ਛਿੱਲੇ}, \emph{maiṁ ālū chhille}.

\begin{cousinweb}
\textbf{\pa{ਛਿੱਲਣਾ}} (\emph{chhillṇā}) is inherited from a reconstructed \emph{*chilla}, \textquotedblleft{}torn, cut\textquotedblright{}; CDIAL lists Punjabi \emph{chillṇā}, \textquotedblleft{}to cut with a knife, peel\textquotedblright{} (CDIAL 5051).
\end{cousinweb}

\begin{grammarlens}[title={What you've built}]
One more word for this chapter.
\end{grammarlens}

\subsection*{Your turn}
\begin{itemize}
\raggedright
  \item \emph{Say it:} \emph{chhillṇā}
  \item \emph{Say it:} \emph{chhillṇā}, once more
  \item \emph{Say it:} \emph{maiṁ chākū nāl ālū chhilldī hāṁ.} — I peel the potatoes with a knife.
  \item \emph{Recall:} say the Punjabi for a pestle, then the Punjabi for to pound, then say \emph{chhillṇā} again
\end{itemize}

\begin{tcolorbox}[breakable,colback=teal!4,colframe=teal!35!black,title={Before you move on}]
What does \pa{ਛਿੱਲਣਾ} mean? (\textquotedblleft{}To peel\textquotedblright{}.) Say it once more.
\end{tcolorbox}
\section[\texorpdfstring{\pa{ਕਤਰਨਾ} (katarnā)}{ਕਤਰਨਾ (katarnā)}]{\pa{ਕਤਰਨਾ} (katarnā) — to chop finely}

\label{lesson:PA-C241-katarna}



\begin{quote}
Before the new one: say the Punjabi for to pound, then the Punjabi for to peel.
\end{quote}

\subsection*{You'll want to know: \pa{ਕਤਰਨਾ}}
\textbf{\pa{ਕਤਰਨਾ}} — \emph{katarnā} — \textquotedblleft{}to chop finely\textquotedblright{}.

\begin{quote}
\pa{ਮੈਂ} \pa{ਸਾਗ} \pa{ਕਤਰਦੀ} \pa{ਹਾਂ।}
\end{quote}

\begin{quote}
\emph{maiṁ sāg katardī hāṁ.}
\end{quote}

\begin{quote}
I chop the greens finely.
\end{quote}

Past (I …): \pa{ਮੈਂ} \pa{ਸਾਗ} \pa{ਕਤਰਿਆ}, \emph{maiṁ sāg katriā}.

\begin{cousinweb}
\textbf{\pa{ਕਤਰਨਾ}} (\emph{katarnā}) is inherited, built on Sanskrit \emph{kartari}, \textquotedblleft{}scissors, knife\textquotedblright{}; CDIAL lists Punjabi \emph{katarnā}, \textquotedblleft{}to clip, shear\textquotedblright{} (CDIAL 2858), and does not give the kitchen sense \textquotedblleft{}chop finely\textquotedblright{}.
\end{cousinweb}

\begin{grammarlens}[title={What you've built}]
One more word for this chapter.
\end{grammarlens}

\subsection*{Your turn}
\begin{itemize}
\raggedright
  \item \emph{Say it:} \emph{katarnā}
  \item \emph{Say it:} \emph{katarnā}, once more
  \item \emph{Say it:} \emph{maiṁ sāg katardī hāṁ.} — I chop the greens finely.
  \item \emph{Recall:} say the Punjabi for to pound, then the Punjabi for to peel, then say \emph{katarnā} again
  \item \emph{Write it:} copy \textbf{\pa{ਕਤਰਨਾ}} once, checking the whole word
\end{itemize}

\begin{tcolorbox}[breakable,colback=teal!4,colframe=teal!35!black,title={Before you move on}]
What does \pa{ਕਤਰਨਾ} mean? (\textquotedblleft{}To chop finely\textquotedblright{}.) Say it once more.
\end{tcolorbox}
\section[\texorpdfstring{\pa{ਚੀਰਨਾ} (chīrnā)}{ਚੀਰਨਾ (chīrnā)}]{\pa{ਚੀਰਨਾ} (chīrnā) — to split; to slit}

\label{lesson:PA-C241-chirna}



\begin{quote}
Before the new one: say the Punjabi for to peel, then the Punjabi for to chop finely.
\end{quote}

\subsection*{You'll want to know: \pa{ਚੀਰਨਾ}}
\textbf{\pa{ਚੀਰਨਾ}} — \emph{chīrnā} — \textquotedblleft{}to split; to slit\textquotedblright{}.

\begin{quote}
\pa{ਮੈਂ} \pa{ਚਾਕੂ} \pa{ਨਾਲ} \pa{ਨਿੰਬੂ} \pa{ਚੀਰਦਾ} \pa{ਹਾਂ।}
\end{quote}

\begin{quote}
\emph{maiṁ chākū nāl nimbū chīrdā hāṁ.}
\end{quote}

\begin{quote}
I split the lemon with a knife.
\end{quote}

Past (I …): \pa{ਮੈਂ} \pa{ਨਿੰਬੂ} \pa{ਚੀਰਿਆ}, \emph{maiṁ nimbū chīriā}.

\begin{cousinweb}
\textbf{\pa{ਚੀਰਨਾ}} (\emph{chīrnā}) is inherited from a reconstructed \emph{*cīrayati}, \textquotedblleft{}splits\textquotedblright{}; CDIAL lists Punjabi \emph{cīrnā}, \textquotedblleft{}to rend, split, saw\textquotedblright{} (CDIAL 4844).
\end{cousinweb}

\begin{grammarlens}[title={What you've built}]
One more word for this chapter.
\end{grammarlens}

\subsection*{Your turn}
\begin{itemize}
\raggedright
  \item \emph{Say it:} \emph{chīrnā}
  \item \emph{Say it:} \emph{chīrnā}, once more
  \item \emph{Say it:} \emph{maiṁ chākū nāl nimbū chīrdā hāṁ.} — I split the lemon with a knife.
  \item \emph{Recall:} say the Punjabi for to peel, then the Punjabi for to chop finely, then say \emph{chīrnā} again
  \item \emph{Write it:} copy \textbf{\pa{ਚੀਰਨਾ}} once, checking the whole word
\end{itemize}

\begin{tcolorbox}[breakable,colback=teal!4,colframe=teal!35!black,title={Before you move on}]
What does \pa{ਚੀਰਨਾ} mean? (\textquotedblleft{}To split; to slit\textquotedblright{}.) Say it once more.
\end{tcolorbox}
\section[\texorpdfstring{\pa{ਮਿਰਚ} (mirach)}{ਮਿਰਚ (mirach)}]{\pa{ਮਿਰਚ} (mirach) — a chilli; pepper}

\label{lesson:PA-C241-mirach}



\begin{quote}
Before the new one: say the Punjabi for to chop finely, then the Punjabi for to split.
\end{quote}

\subsection*{You'll want to know: \pa{ਮਿਰਚ}}
\textbf{\pa{ਮਿਰਚ}} — \emph{mirach} — \textquotedblleft{}a chilli; pepper\textquotedblright{}.

\begin{quote}
\pa{ਇਹ} \pa{ਹਰੀ} \pa{ਮਿਰਚ} \pa{ਬਹੁਤ} \pa{ਤਿੱਖੀ} \pa{ਹੈ।}
\end{quote}

\begin{quote}
\emph{eh harī mirach bahut tikkhī hai.}
\end{quote}

\begin{quote}
This green chilli is very hot.
\end{quote}

A feminine noun.

\begin{cousinweb}
\textbf{\pa{ਮਿਰਚ}} (\emph{mirach}) is inherited from Sanskrit \emph{marīca}, \textquotedblleft{}peppercorn\textquotedblright{}; CDIAL lists Punjabi \emph{mirc}, \textquotedblleft{}black pepper\textquotedblright{} (CDIAL 9875). CDIAL does not give the everyday sense \textquotedblleft{}chilli\textquotedblright{}, and it notes that the Sanskrit word may itself be Austroasiatic.
\end{cousinweb}

\begin{grammarlens}[title={What you've built}]
One more word for this chapter.
\end{grammarlens}

\subsection*{Your turn}
\begin{itemize}
\raggedright
  \item \emph{Say it:} \emph{mirach}
  \item \emph{Say it:} \emph{mirach}, once more
  \item \emph{Say it:} \emph{eh harī mirach bahut tikkhī hai.} — This green chilli is very hot.
  \item \emph{Recall:} say the Punjabi for to chop finely, then the Punjabi for to split, then say \emph{mirach} again
  \item \emph{Write it:} copy \textbf{\pa{ਮਿਰਚ}} once, checking the whole word
\end{itemize}

\begin{tcolorbox}[breakable,colback=teal!4,colframe=teal!35!black,title={Before you move on}]
What does \pa{ਮਿਰਚ} mean? (\textquotedblleft{}A chilli; pepper\textquotedblright{}.) Say it once more.
\end{tcolorbox}
\section[\texorpdfstring{\pa{ਵੇਸਣ} (vesaṇ)}{ਵੇਸਣ (vesaṇ)}]{\pa{ਵੇਸਣ} (vesaṇ) — gram flour}

\label{lesson:PA-C241-vesan}



\begin{quote}
Before the new one: say the Punjabi for to split, then the Punjabi for a chilli.
\end{quote}

\subsection*{You'll want to know: \pa{ਵੇਸਣ}}
\textbf{\pa{ਵੇਸਣ}} — \emph{vesaṇ} — \textquotedblleft{}gram flour\textquotedblright{}.

\begin{quote}
\pa{ਮੈਂ} \pa{ਵੇਸਣ} \pa{ਨਾਲ} \pa{ਪਕੌੜੇ} \pa{ਬਣਾਉਂਦੀ} \pa{ਹਾਂ।}
\end{quote}

\begin{quote}
\emph{maiṁ vesaṇ nāl pakauṛe baṇāundī hāṁ.}
\end{quote}

\begin{quote}
I make pakoras with gram flour.
\end{quote}

A masculine noun.

\begin{cousinweb}
\textbf{\pa{ਵੇਸਣ}} (\emph{vesaṇ}) is inherited from Sanskrit \emph{vēsana}, \textquotedblleft{}flour of a particular vegetable product\textquotedblright{}; CDIAL lists Punjabi \emph{vesaṇ}, \textquotedblleft{}gram flour\textquotedblright{} (CDIAL 12133).
\end{cousinweb}

\begin{grammarlens}[title={What you've built}]
One more word for this chapter.
\end{grammarlens}

\subsection*{Your turn}
\begin{itemize}
\raggedright
  \item \emph{Say it:} \emph{vesaṇ}
  \item \emph{Say it:} \emph{vesaṇ}, once more
  \item \emph{Say it:} \emph{maiṁ vesaṇ nāl pakauṛe baṇāundī hāṁ.} — I make pakoras with gram flour.
  \item \emph{Recall:} say the Punjabi for to split, then the Punjabi for a chilli, then say \emph{vesaṇ} again
  \item \emph{Write it:} copy \textbf{\pa{ਵੇਸਣ}} once, checking the whole word
\end{itemize}

\begin{tcolorbox}[breakable,colback=teal!4,colframe=teal!35!black,title={Before you move on}]
What does \pa{ਵੇਸਣ} mean? (\textquotedblleft{}Gram flour\textquotedblright{}.) Say it once more.
\end{tcolorbox}
